% GENERATED FILE. Edit canonical lessons, then run npm run generate:books.
% canonical-source-hash: fnv1a64:389016e0773c8623
% canonical-lessons: FR-C282-gros-titre, FR-C282-manchette, FR-C282-rubrique, FR-C282-editorial, FR-C282-faits-divers

\chapter{Les gros titres}
\label{ch:fr-les-gros-titres}
\hlchaptermodality{282}

\begin{chapteropening}
\textbf{By the end of this chapter:} \emph{I can talk about the front page: headlines, banner headlines, sections, editorials and short news items.}

Complete the last lesson of chapter 282: I can talk about the front page: headlines, banner headlines, sections, editorials and short news items.
\end{chapteropening}

\section[le gros titre]{le gros titre — the headline}

\label{lesson:FR-C282-gros-titre}



\begin{quote}
Before the new one: say the French for the dispatch, then the French for the news agency.
\end{quote}

\subsection*{You'll want to know: le gros titre}
\textbf{le gros titre} — \textquotedblleft{}the headline\textquotedblright{}.

\begin{quote}
Le gros titre du jour parle du prix du pain.
\end{quote}

\begin{quote}
\emph{Today's headline is about the price of bread.}
\end{quote}

\emph{Faire les gros titres} — to make the headlines. \emph{À la une} — on the front page.

\begin{cousinweb}
\emph{titre} is inherited from Latin \emph{titulus} \textquotedblleft{}inscription, label\textquotedblright{} (EtymDB), and \emph{gros} from Late Latin \emph{grossus} \textquotedblleft{}thick\textquotedblright{} (EtymDB).
\end{cousinweb}

\begin{grammarlens}[title={What you've built}]
One more word for telling, explaining and describing at B1.
\end{grammarlens}

\subsection*{Your turn}
\begin{itemize}
\raggedright
  \item \emph{Say it:} \emph{le gros titre}
  \item \emph{Say it:} \emph{le gros titre}, once more
  \item \emph{Say it:} a sentence of your own with \emph{le gros titre}
  \item \emph{Recall:} say the French for the dispatch, then the French for the news agency, then say \emph{le gros titre} again
\end{itemize}

\begin{tcolorbox}[breakable,colback=teal!4,colframe=teal!35!black,title={Before you move on}]
What does le gros titre mean? (\textquotedblleft{}The headline\textquotedblright{}.) Say it once more.
\end{tcolorbox}
\section[la manchette]{la manchette — the banner headline; a shirt cuff}

\label{lesson:FR-C282-manchette}



\begin{quote}
Before the new one: say the French for the news agency, then the French for the headline.
\end{quote}

\subsection*{You'll want to know: la manchette}
\textbf{la manchette} — \textquotedblleft{}the banner headline; a shirt cuff\textquotedblright{}.

\begin{quote}
La manchette annonçait la victoire de l'équipe.
\end{quote}

\begin{quote}
\emph{The banner headline announced the team's victory.}
\end{quote}

On a shirt, \emph{les manchettes} are the cuffs.

\begin{cousinweb}
\emph{manchette} is a diminutive of \emph{manche} \textquotedblleft{}sleeve\textquotedblright{} (EtymDB), which is inherited from Latin \emph{manica} \textquotedblleft{}a long sleeve\textquotedblright{}, built on \emph{manus} \textquotedblleft{}hand\textquotedblright{} (Lewis \& Short).
\end{cousinweb}

\begin{grammarlens}[title={What you've built}]
One more word for telling, explaining and describing at B1.
\end{grammarlens}

\subsection*{Your turn}
\begin{itemize}
\raggedright
  \item \emph{Say it:} \emph{la manchette}
  \item \emph{Say it:} \emph{la manchette}, once more
  \item \emph{Say it:} a sentence of your own with \emph{la manchette}
  \item \emph{Recall:} say the French for the news agency, then the French for the headline, then say \emph{la manchette} again
\end{itemize}

\begin{tcolorbox}[breakable,colback=teal!4,colframe=teal!35!black,title={Before you move on}]
What does la manchette mean? (\textquotedblleft{}The banner headline; a shirt cuff\textquotedblright{}.) Say it once more.
\end{tcolorbox}
\section[la rubrique]{la rubrique — the section, the regular feature (in a paper)}

\label{lesson:FR-C282-rubrique}



\begin{quote}
Before the new one: say the French for the headline, then the French for the banner headline.
\end{quote}

\subsection*{You'll want to know: la rubrique}
\textbf{la rubrique} — \textquotedblleft{}the section, the regular feature (in a paper)\textquotedblright{}.

\begin{quote}
Il écrit dans la rubrique sport.
\end{quote}

\begin{quote}
\emph{He writes in the sports section.}
\end{quote}

\emph{La rubrique nécrologique} — the death notices.

\begin{cousinweb}
\emph{rubrique} is a learned borrowing from Latin \emph{rubrīca} \textquotedblleft{}red earth, a heading written in red\textquotedblright{}, from \emph{ruber} \textquotedblleft{}red\textquotedblright{} (Lewis \& Short).
\end{cousinweb}

\begin{grammarlens}[title={What you've built}]
One more word for telling, explaining and describing at B1.
\end{grammarlens}

\subsection*{Your turn}
\begin{itemize}
\raggedright
  \item \emph{Say it:} \emph{la rubrique}
  \item \emph{Say it:} \emph{la rubrique}, once more
  \item \emph{Say it:} a sentence of your own with \emph{la rubrique}
  \item \emph{Recall:} say the French for the headline, then the French for the banner headline, then say \emph{la rubrique} again
\end{itemize}

\begin{tcolorbox}[breakable,colback=teal!4,colframe=teal!35!black,title={Before you move on}]
What does la rubrique mean? (\textquotedblleft{}The section, the regular feature (in a paper)\textquotedblright{}.) Say it once more.
\end{tcolorbox}
\section[l'éditorial]{l'éditorial — the editorial, the leading article}

\label{lesson:FR-C282-editorial}



\begin{quote}
Before the new one: say the French for the banner headline, then the French for the section.
\end{quote}

\subsection*{You'll want to know: l'éditorial}
\textbf{l'éditorial} — \textquotedblleft{}the editorial, the leading article\textquotedblright{}.

\begin{quote}
Le directeur du journal écrit l'éditorial.
\end{quote}

\begin{quote}
\emph{The editor of the paper writes the editorial.}
\end{quote}

Plural \emph{les éditoriaux}. Its writer is \emph{l'éditorialiste}.

\begin{cousinweb}
\emph{éditorial} is usually traced to English \emph{editorial} (EtymDB lists no parent), built on \emph{editor}, which comes from Latin \emph{ēditor} \textquotedblleft{}one who brings out\textquotedblright{}, from \emph{ēdere} \textquotedblleft{}to give out\textquotedblright{} (Lewis \& Short).
\end{cousinweb}

\begin{grammarlens}[title={What you've built}]
One more word for telling, explaining and describing at B1.
\end{grammarlens}

\subsection*{Your turn}
\begin{itemize}
\raggedright
  \item \emph{Say it:} \emph{l'éditorial}
  \item \emph{Say it:} \emph{l'éditorial}, once more
  \item \emph{Say it:} a sentence of your own with \emph{l'éditorial}
  \item \emph{Recall:} say the French for the banner headline, then the French for the section, then say \emph{l'éditorial} again
\end{itemize}

\begin{tcolorbox}[breakable,colback=teal!4,colframe=teal!35!black,title={Before you move on}]
What does l'éditorial mean? (\textquotedblleft{}The editorial, the leading article\textquotedblright{}.) Say it once more.
\end{tcolorbox}
\section[les faits divers]{les faits divers — short news items, minor news stories}

\label{lesson:FR-C282-faits-divers}



\begin{quote}
Before the new one: say the French for the section, then the French for the editorial.
\end{quote}

\subsection*{You'll want to know: les faits divers}
\textbf{les faits divers} — \textquotedblleft{}short news items, minor news stories\textquotedblright{}.

\begin{quote}
Il aime lire les faits divers.
\end{quote}

\begin{quote}
\emph{He likes reading the short news items.}
\end{quote}

Small stories: thefts, odd events, accidents. Singular: \emph{un fait divers}.

\begin{cousinweb}
\emph{fait} is inherited from Latin \emph{factum} \textquotedblleft{}a deed\textquotedblright{} (Lewis \& Short), and \emph{divers} is a learned borrowing from Latin \emph{dīversus} \textquotedblleft{}turned different ways\textquotedblright{} (EtymDB).
\end{cousinweb}

\begin{grammarlens}[title={What you've built}]
One more word for telling, explaining and describing at B1.
\end{grammarlens}

\subsection*{Your turn}
\begin{itemize}
\raggedright
  \item \emph{Say it:} \emph{les faits divers}
  \item \emph{Say it:} \emph{les faits divers}, once more
  \item \emph{Say it:} a sentence of your own with \emph{les faits divers}
  \item \emph{Recall:} say the French for the section, then the French for the editorial, then say \emph{les faits divers} again
\end{itemize}

\begin{tcolorbox}[breakable,colback=teal!4,colframe=teal!35!black,title={Before you move on}]
What does les faits divers mean? (\textquotedblleft{}Short news items, minor news stories\textquotedblright{}.) Say it once more.
\end{tcolorbox}
